\section{Билет 8. Виды колебаний. Осциллятор. Потенциальная энергия вблизи положения равновесия. Квазиупругая сила. Гармонический осциллятор и его примеры.}

\begin{definition}
В зависимости от характера воздействия на систему различают следующие основные виды колебаний:
\begin{enumerate}
    \item \textbf{Свободные (собственные) колебания} --- происходят в системе, предоставленной самой себе, после выведения из положения равновесия или сообщения начального импульса (внешние силы отсутствуют).
    \item \textbf{Вынужденные колебания} --- протекают под действием внешней периодически изменяющейся силы.
    \item \textbf{Автоколебания} --- незатухающие колебания, поддерживаемые постоянным источником энергии за счёт самой колебательной системы (например, маятниковые часы с анкерным ходом).
    \item \textbf{Параметрические колебания} --- возникают при периодическом изменении какого-либо параметра системы (например, длины маятника, индуктивности катушки или ёмкости конденсатора).
\end{enumerate}
\end{definition}

\begin{definition}
\textbf{Осциллятором} называют систему, находящуюся в положении устойчивого равновесия, которая при небольших отклонениях стремится вернуться обратно. Устойчивое равновесие соответствует локальному минимуму потенциальной энергии $U(x)$.
\end{definition}

\begin{theorem}[Потенциальная энергия вблизи положения равновесия]
Рассмотрим одномерное движение в поле с потенциальной энергией $U(x)$. Положим начало координат в точку устойчивого равновесия ($x=0$) и отсчитываем энергию от минимума: $U(0)=0$. Разложим $U(x)$ в ряд Тейлора по степеням $x$:
\begin{equation}
    U(x) = U(0) + U'(0)x + \frac{1}{2}U''(0)x^2 + \frac{1}{6}U'''(0)x^3 + \dots
\end{equation}
Так как в точке $x=0$ потенциал имеет минимум, то первая производная равна нулю: $U'(0)=0$, а вторая производная положительна: $U''(0) = k > 0$. Пренебрегая членами высшего порядка малости ($x^3, x^4, \dots$) при малых отклонениях, получаем квадратичное приближение:
\begin{equation}
    U(x) \approx \frac{1}{2}kx^2. \label{eq:potential_approx}
\end{equation}
\end{theorem}

\begin{definition}
Сила, связанная с потенциальной энергией соотношением $F_x = -\frac{dU}{dx}$, вблизи положения равновесия принимает вид:
\begin{equation}
    F_x = -kx. \label{eq:quasi_elastic}
\end{equation}
Знак «минус» указывает на то, что сила всегда направлена к положению равновесия ($x=0$). Силы вида $F = -kx$, независимо от их физической природы (гравитационные, электрические, упругие и т.д.), называются \textbf{квазиупругими}. Коэффициент $k$ в данном случае является эффективным коэффициентом жёсткости системы.
\end{definition}

\begin{theorem}[Гармонический осциллятор]
Подставляя квазиупругую силу \eqref{eq:quasi_elastic} во второй закон Ньютона $m\ddot{x} = F_x$, получаем уравнение движения:
\begin{equation}
    m\ddot{x} = -kx \quad \Rightarrow \quad \ddot{x} + \frac{k}{m}x = 0.
\end{equation}
Вводя $\omega_0^2 = k/m$, приходим к каноническому уравнению:
\begin{equation}
    \ddot{x} + \omega_0^2 x = 0. \label{eq:harmonic_osc}
\end{equation}
Система, описываемая данным уравнением, называется \textbf{гармоническим осциллятором}. Это идеализированная модель, совершающая свободные незатухающие гармонические колебания около положения устойчивого равновесия.
\end{theorem}

\begin{remark}
\textbf{Примеры гармонического осциллятора в физике:}
\begin{itemize}
    \item Математический и физический маятники (в приближении малых углов);
    \item Груз на пружине (пружинный маятник);
    \item Крутильный маятник;
    \item Колебательный электрический контур (LC-контур);
    \item Колебания атомов в кристаллической решётке твёрдых тел и молекулах.
\end{itemize}
Несмотря на различную физическую природу, все эти системы в малом приближении описываются одним и тем же математическим аппаратом гармонического осциллятора.
\end{remark}
